% Copyright 2004 by Till Tantau <tantau@users.sourceforge.net>.
%
% In principle, this file can be redistributed and/or modified under
% the terms of the GNU Public License, version 2.
%
% However, this file is supposed to be a template to be modified
% for your own needs. For this reason, if you use this file as a
% template and not specifically distribute it as part of a another
% package/program, I grant the extra permission to freely copy and
% modify this file as you see fit and even to delete this copyright
% notice. 

\documentclass{beamer}
\usefonttheme[onlylarge]{structurebold}
\setbeamerfont*{frametitle}{size=\normalsize,series=\bfseries}
\setbeamertemplate{navigation symbols}{}
% Altera as cores do template
%\setbeamercolor{title}{fg=cyan!30!darkgray}
%\setbeamercolor{institute}{fg=cyan!30!darkgray}
%\setbeamercolor{author}{fg=cyan!30!darkgray}
%\setbeamercolor{date}{fg=cyan!30!darkgray}
%\setbeamercolor{frametitle}{fg=cyan!30!darkgray}
%\setbeamercolor{section in toc}{fg=cyan!30!darkgray}

%\setbeamercolor{section in sidebar}{fg=red}
%\setbeamercolor{subsection in sidebar}{bg=cyan!30!darkgray}
%\setbeamercolor{palette sidebar tertiary}{fg=cyan!20!gray}
%\setbeamercolor{palette sidebar quaternary}{fg=cyan!20!gray}

%\setbeamercolor{item}{fg=red}



% Standard packages

\usepackage[brazil]{babel}
\usepackage[latin1]{inputenc}
\usepackage{times}
\usepackage[T1]{fontenc}
\usepackage{multicol}
\usepackage{amssymb}
\usepackage{hyperref}
\usepackage{textcomp}
\usepackage{graphics}
\usepackage{graphicx}
\usepackage{ragged2e} % Par�grafo justificado
%\usepackage[table]{xcolor}
 
% Setup TikZ

\usepackage{tikz}
\usetikzlibrary{arrows}
\tikzstyle{block}=[draw opacity=0.7,line width=1.4cm]

%diretório das figuras
%\graphicspath{./pictures}


% There are many different themes available for Beamer. A comprehensive
% list with examples is given here:
% http://deic.uab.es/~iblanes/beamer_gallery/index_by_theme.html
% You can uncomment the themes below if you would like to use a different
% one:
%\usetheme{AnnArbor}
%\usetheme{Antibes}
%\usetheme{Bergen}
%\usetheme{Berkeley}
%\usetheme{Berlin}
%\usetheme{Boadilla}
%\usetheme{boxes}
%\usetheme{CambridgeUS}
%\usetheme{Copenhagen}
%\usetheme{Darmstadt}
%\usetheme{default}
%\usetheme{Frankfurt}
\usetheme{Goettingen}
%\usetheme{Hannover}
%\usetheme{Ilmenau}
%\usetheme{JuanLesPins}
%\usetheme{Luebeck}
%\usetheme{Madrid}
%\usetheme{Malmoe}
%\usetheme{Marburg}
%\usetheme{Montpellier}
%\usetheme{PaloAlto}
%\usetheme{Pittsburgh}
%\usetheme{Rochester}
%\usetheme{Singapore}
%\usetheme{Szeged}
%\usetheme{Warsaw}




\title{IPTables}

% A subtitle is optional and this may be deleted
%\subtitle{Optional Subtitle}

\author{S. Danilo\inst{1}}
% - Give the names in the same order as the appear in the paper.
% - Use the \inst{?} command only if the authors have different
%   affiliation.

\institute[Universidade Federal do Para] % (optional, but mostly needed)
{
  \inst{1}%
  Divis�o de Servi�os de Internet \\
  Coordenadoria de Seguran�a e  Servi�os de Internet \\
  Datacenter da UFPA \\
% - Use the \inst command only if there are several affiliations.
% - Keep it simple, no one is interested in your street address.
}

\date{Bel�m, 2013}
% - Either use conference name or its abbreviation.
% - Not really informative to the audience, more for people (including
%   yourself) who are reading the slides online

\subject{IPTables}
% This is only inserted into the PDF information catalog. Can be left
% out. 

% If you have a file called "university-logo-filename.xxx", where xxx
% is a graphic format that can be processed by latex or pdflatex,
% resp., then you can add a logo as follows:

% \pgfdeclareimage[height=0.5cm]{university-logo}{university-logo-filename}
% \logo{\pgfuseimage{university-logo}}

% Delete this, if you do not want the table of contents to pop up at
% the beginning of each subsection:
%\AtBeginSubsection[]
%{
%  \begin{frame}<beamer>{Outline}
%    \tableofcontents[currentsection,currentsubsection]
%  \end{frame}
%}

% Let's get started
\begin{document}

\begin{frame}
	\begin{figure}[htb]
		\centering 
		\includegraphics[scale=0.6]{./pictures/ctic.png}
		\label{fig:ctic}			
	\end{figure}
  	\titlepage
\end{frame}

\begin{frame}{Agenda}
  \tableofcontents
  % You might wish to add the option [pausesections]
\end{frame}

\section{Introdu��o � Firewall}

	\begin{frame}
		\begin{block}{Defini��o}
			Um conjunto de aplica��es que realiza controle de entrada e sa�da de uma determinada rede, sobre qualquer uma das camadas da pilha 					TCP/IP
		\end{block}
		\begin{itemize}
			\item \textit{NetFilter} - Pacotes
				\begin{itemize}
					\item \textit{Stateless}
						\begin{itemize}
							\item Toma decis�es baseado somente no pacote atual
							\item Muito utilizado em roteadores simples
							\item Baixo custo computacional
						\end{itemize}
					\item \textit{Statefull}
						\begin{itemize}
							\item Toma decis�es baseado em informa�es armazenadas
							\item Mais utilizado em servidores
							\item Alto custo computacional
						\end{itemize}			
				\end{itemize}
			\item Squid (Proxy)
				\begin{itemize}
					\item Filtro de aplica��o
					\item Bloqueio de url\textquotesingle s e palavras chave
				\end{itemize} 
		\end{itemize}
		
		\begin{alertblock}{Importante}
			Um firewall mal configurado � pior do que n�o ter um firewall	
		\end{alertblock}
	\end{frame}
	
	\begin{frame}{Prote��o}
		\begin{itemize}
			\item Firewall de \textit{host}
				\begin{itemize}
					\item Protege somente a m�quina local
					\item Regras mais robustas
				\end{itemize}
			\item Firewall de rede
				\begin{itemize}
					\item Protege a(s) rede(s) � qual pertence
					\item Regras mais simples
					\item \alert {Pode se tornar um gargalo}
				\end{itemize}
		\end{itemize}
	\end{frame}

\section{IPTables}

\subsection{Introdu��o}
	\begin{frame}{Introdu��o}
		\begin{itemize}
			\item A partir do kernel 2.4
			\item Manipula o \textit{netfilter} (usu�rio)
			\item Est� implementado diretamente no kernel
			\item Firewall \textit{Statefull}
			\item Firewall de \textit{host}/rede
		\end{itemize}
	\end{frame}

\subsection{Estrutura}
	\begin{frame}
		
	\end{frame}

\subsubsection{Tabela Filter}

\subsubsection{Tabela Nat}

\section{Exemplos}
						
\end{document}
